\documentclass[10pt,a4paper]{amsart}
\usepackage[margin=19mm]{geometry}
\usepackage{amsmath,amssymb,mathtools}
\usepackage[scr=rsfs,cal=euler]{mathalfa}
\usepackage{xfrac}
\usepackage[unicode,hidelinks,bookmarks=false]{hyperref}
\newcommand{\ie}{i.e.\ }
\newcommand{\cf}{cf.\ }
\newcommand{\resp}{respectively\ }
\newcommand{\Subsection}{\subsection}
\providecommand{\nobreakdash}{\nobreak\hbox{-}}
\setlength{\emergencystretch}{2em}
\multlinegap0pt
\usepackage{bm}
\usepackage{bbm}
\usepackage{dsfont}
\usepackage{xspace}
\usepackage{cancel}
\usepackage{mathtools}
\usepackage{amsthm}
\makeatletter
\@ifundefined{theorem}{\newtheorem{theorem}{Theorem}}{}
\@ifundefined{lemma}{\newtheorem{lemma}[theorem]{Lemma}}{}
\makeatother
\newcommand{\Q}{\mathbb{Q}}
\renewcommand{\le}{\leqslant}
\title[The non-existence of some Galois representations of moderate]{The non-existence of some Galois representations of moderate dimension in small characteristic}
\author{Alexandru Ghitza, Takuya Yamauchi}
\date{Source: arXiv:2509.00635v2 (CC BY 4.0)}
\begin{document}
\maketitle

\noindent\emph{Source and licence.} The non-existence of some Galois representations of moderate dimension in small characteristic. Version arXiv:2509.00635v2 (CC BY 4.0), distributed under the Creative Commons Attribution 4.0 International licence, as recorded in the arXiv licence metadata for this exact version and shown on the abstract page https://arxiv.org/abs/2509.00635v2. Full rights notice: licenses/arxiv-2509.00635-CC-BY-4.0.txt.\par\smallskip

\noindent\textit{Editorial scope.} Counted unit: the Lemma whose source label is \texttt{poly} in the article (source0.tex lines 673--677, source label \texttt{poly}), delivered together with the article's own complete original proof of it (source0.tex lines 678--721, one \texttt{proof} environment of 44 lines). No mathematics is written, repaired or re-derived: statement and proof are the article's own text line for line. Required context retained uncounted: none. Every definition, display and label of those lines is retained unchanged. Omitted from this package and not redistributed: 1--672, 722--1259. Nothing inside a retained excerpt is omitted or altered: every retained body is delivered exactly as the article's lines are written, so no omission receipt of any kind accompanies this package. Citations: the retained excerpts carry no citation command at all, so no bibliography entry of the article is transcribed and no reference is redistributed. Reference adaptation: none was needed and none was made; every reference inside the delivered unit resolves inside it, so no unresolved reference or citation is produced. Printed slips kept literal: no printed word, symbol, hypothesis or number of the counted statement or of its proof was altered. Layout and class: the article's own class is \texttt{amsart} with packages including amsmath, amssymb, amscd, latexsym, graphicx, mathrsfs, enumerate, mathtools, color, booktabs, hyperref; the wrapper is the lane's pinned \texttt{amsart} editorial wrapper at 10pt on A4, which restates the theorem-like environments the retained text needs and adds the article's own macro definitions that the retained text needs (\texttt{\textbackslash Q}, \texttt{\textbackslash le}), with extra wrapper packages (bm, bbm, dsfont, xspace, cancel, mathtools). The delivered numbering is the unit's own. Assets: no retained span contains a figure environment, a graphics-inclusion command, a PDF-page inclusion, a table or a TikZ picture; no image file is copied into this package, the package declares no asset dependency and no image file is redistributed with it. Licence: version arXiv:2509.00635v2 of this article is distributed under the Creative Commons Attribution 4.0 International licence, as recorded in the arXiv licence metadata for this exact version and shown on the abstract page https://arxiv.org/abs/2509.00635v2. A full rights notice is delivered at licenses/arxiv-2509.00635-CC-BY-4.0.txt. Characters: every retained excerpt is pure ASCII, so the delivered engine, XeLaTeX with its default Latin Modern text font, reports no missing character.
\begin{lemma}\label{poly}Let $K/\Q$ be a Galois extension with Galois group $G$
  and an injective group homomorphism $\iota:G\hookrightarrow S_n$.
  There exists a separable polynomial $f(x)\in\Q[x]$ of degree at most $n$
  such that $K=\Q_f$, the splitting field of $f$. Further, if $G$ acts transitively on $\{1,2,\ldots,n\}$, then $f(x)$ is irreducible over $\Q$ of degree $n$.
\end{lemma}

\begin{proof}Put $H:=\iota(G)\subset S_n$.
  For each $i$ with $1\le i\le n$, let $H_i$ be the stabilizer of $i$ and put
  $G_i=\iota^{-1}(H_i)\subset G$. We also define
  \begin{equation*}
    N_i:=\bigcap_{\sigma\in G}\sigma G_i \sigma^{-1}.
  \end{equation*}
  Let us consider the orbit decomposition
  \begin{equation*}
    \{1,2,\ldots,n\}=\coprod_{k=1}^r O_H(i_k).
  \end{equation*}
  For each $i_k$,
  put $K_{i_k}:=K^{G_{i_k}}$ and choose $\alpha_{i_k}\in K$ such that
  $K_{i_k}=K^{G_{i_k}}=\Q(\alpha_{i_k})$. For each $\sigma \in G$, $\sigma(K_{i_k})=
    K^{\sigma G_{i_k}\sigma^{-1}}$, so the composite of all conjugates of $K_{i_k}$
  is $K^{N_{i_k}}$.
  Let $f_{i_k}(x)$ be the minimal polynomial of $\alpha_{i_k}$ and define
  \begin{equation}\label{g(x)}
    \widetilde{f}(x):=\prod_{k=1}^r f_{i_k}(x).
  \end{equation}
  It is easy to see that
  \begin{equation*}
    \deg(\widetilde{f}(x))=\sum_{k=1}^r[K_{i_k}:\Q]=\sum_{k=1}^r[G:G_{i_k}]=\sum_{k=1}^r[H:H_{i_k}]=
    \sum_{k=1}^r|O_H(i_k)|=n.
  \end{equation*}
  We show that $\Q_{\widetilde{f}}=K$: since $\iota$ is injective, we have
  \begin{equation*}
    \bigcap_{k=1}^r N_{i_k}\subset \bigcap_{j=1}^n G_{j}=\iota^{-1}(\bigcap_{j=1}^n H_j)
    =\iota^{-1}(\bigcap_{k=1}^r H_{i_k})=\iota^{-1}(\{1\})=\{1\},
  \end{equation*}
  therefore
  \begin{equation*}
    \Q_{\widetilde{f}}=\prod_{k=1}^r K^{N_{i_k}}=K^{\bigcap_{k=1}^r N_{i_k}}=K.
  \end{equation*}
  It remains to deal with the (possible) non-separability of $\widetilde{f}$.
  If there exist $k,k'\in\{1,\dots,r\}$ such that $\alpha_{i_k}$ is conjugate under the Galois action of $G$ to $\alpha_{i_{k'}}$, then $f_{i_k}(x)=f_{i_{k'}}(x)$, so we can keep $f_{i_k}(x)$ and discard $f_{i_{k'}(x)}$.
  We obtain a subset $S$ of $\{1,\ldots, r\}$ such that
  \begin{equation*}
    f(x):=\prod_{k\in S}f_{i_k}(x)
  \end{equation*}
  is separable and $\Q_f=\Q_{\widetilde{f}}=K$. Clearly
  $\deg(f(x))\le \deg(\widetilde{f}(x))=n$.

  The latter claim follows easily from the construction.
\end{proof}

\end{document}
